% !TeX root = ../main.tex
\subsection{تمرین \secnum{۴.۷}، بازاندازه‌گذاری و تعداد مراحل}
\label{sec:delay-exercise-stage-count}

\textbf{خواسته.}
برای مسیر سه‌مرحله‌ای با \(f_1=12\)، \(f_2=6\) و \(f_3=9\)
بررسی کنید آیا می‌توان تأخیر را کاهش داد. تعداد مراحل مناسب و
تغییرهای پیشنهادی را تعیین کنید
\cite[تمرین \secnum{۴.۷}، صفحه ۱۷۷]{weste}.

\textbf{بازاندازه‌گذاری مسیر موجود.}
فرض می‌کنیم توپولوژی، بار انتهایی، ظرفیت ورودی مسیر و ضرایب
انشعاب ثابت‌اند و اندازه‌های میانی را می‌توان پیوسته تغییر داد.
حاصل‌ضرب \(F=12\cdot6\cdot9=648\) و مجموع اولیه
\(12+6+9=27\) است.
برای \(F=\prod_i f_i\) ثابت، مجموع \(f_i\)ها زمانی کمینه است
که همه برابر با میانگین هندسی باشند. بنابراین مقدار مشترکِ بهینه \(\widehat f=\sqrt[3]{648}\approx8.65\)
و مجموع جدید تقریباً \num{25.96} است. اگر \(P_0\) مجموع
تأخیر \lr{Parasitic} سه گیت اولیه باشد، تأخیر در مدل از
\(27+P_0\) به حدود \(25.96+P_0\) می‌رسد.

برای دیدن جهت تغییر اندازه‌ها، \(C_2\) و \(C_3\) را ظرفیت
ورودی مراحل دوم و سوم در مدار اولیه بگیرید و اندازه‌های
بازتنظیم‌شده را با ستاره مشخص کنید. با \(g_1,b_1,C_1\)
ثابت، \(f_1\) با \(C_2\) متناسب است. از دید مرحله سوم،
بار نهایی ثابت است و \(f_3\) با \(C_3\) نسبت معکوس دارد.
در نتیجه، برای برابرشدن هر سه \lr{Stage Effort} با
\(\widehat f\) داریم.
\[
  \frac{C_2^\star}{C_2}=\frac{\widehat f}{12}\approx0.72,
  \qquad
  \frac{C_3^\star}{C_3}=\frac{9}{\widehat f}\approx1.04.
\]
ثابت‌ماندن \(F\) تضمین می‌کند که پس از این تغییر، \(f_2\)
نیز به \(\widehat f\) برسد. پس ظرفیت ورودی مرحله دوم
کاهش و ظرفیت ورودی مرحله سوم اندکی افزایش می‌یابد. اگر
انشعاب‌ها با این تغییر رفتار متفاوتی داشته باشند، باید
ضرایب \(b_i\) را دوباره حساب کرد.

\textbf{مقایسه تعداد مراحل.}
اگر گیت‌های سه مرحله اول حفظ شوند و بتوان به انتهای مسیر
\lr{Inverter} اضافه کرد،
در مدل با بار انتهایی و ضریب‌های انشعاب ثابت، \(F=648\) ثابت است و
هر مرحله تازه یک واحد تأخیر بی‌بار اضافه می‌کند.
\(\log_4 648\approx4.67\) نقطه شروع مقایسه است.
در تقریب پیوسته، \(\rho=F^{1/N}\) را تعریف می‌کنیم.
مشتق \(NF^{1/N}+(N-3)p_{\mathrm{inv}}\) نسبت به \(N\)
برابر \(p_{\mathrm{inv}}+\rho(1-\ln\rho)\) است.
صفرکردن مشتق با \(p_{\mathrm{inv}}=1\)، حدود \num{5.06} مرحله می‌دهد.
جدول \ref{tab:delay-exercise-stage-count} مقادیر
\(D_{\min}(N)-P_0\) را از
محاسبه \(NF^{1/N}+(N-3)p_{\mathrm{inv}}\) نشان می‌دهد.

\begin{table}[htbp]
  \centering
  \caption{مقایسه تعداد مراحل در تمرین \secnum{۴.۷}.
  \(D_{\min}\) و \(P_0\) نرمال‌شده‌اند و مقادیر ستون تأخیر
  تا دو رقم اعشار گرد شده‌اند.}
  \label{tab:delay-exercise-stage-count}
  \begin{tabular}{cc}
    \toprule
    تعداد مراحل \(N\) & \(D_{\min}(N)-P_0\) \\
    \midrule
    \num{3} & \num{25.96} \\
    \num{4} & \num{21.18} \\
    \num{5} & \num{20.25} \\
    \num{6} & \num{20.65} \\
    \num{7} & \num{21.65} \\
    \bottomrule
  \end{tabular}
\end{table}

در این خانواده از طرح‌ها، پنج مرحله کمترین تأخیر را دارد.
افزودن دو \lr{Inverter} قطبیت اصلی خروجی را حفظ می‌کند.
برای مسیر پنج‌مرحله‌ای، \(\widehat f=\sqrt[5]{648}\approx3.65\)
است و اندازه‌ها را باید از خروجی به ورودی تنظیم کرد.
سه مقدار داده‌شده برای تعیین بهترین توپولوژی میان همه پیاده‌سازی‌های
ممکن کافی نیستند، زیرا با تغییر نوع گیت‌ها یا انشعاب‌ها،
\(F\) و مجموع تأخیر بی‌بار می‌توانند تغییر کنند.
